% Content companion to plastic_plane_strain_case.pdf.
% The supplied PDF is rendered by generate_plastic_plane_strain_case.py (Matplotlib).
% This self-contained source may also be compiled with pdflatex.
\documentclass[10pt,a4paper]{article}
\usepackage[margin=14.8mm]{geometry}
\usepackage[T1]{fontenc}
\usepackage{lmodern}
\usepackage{amsmath,amssymb,array,xcolor}
\usepackage{microtype}
\definecolor{accent}{HTML}{176475}
\renewcommand{\familydefault}{\sfdefault}
\pagestyle{empty}
\setlength{\parindent}{0pt}
\setlength{\parskip}{3pt}
\setlength{\abovedisplayskip}{4pt}
\setlength{\belowdisplayskip}{4pt}
\newcommand{\topic}[1]{\par\vspace{4pt}{\color{accent}\bfseries\large #1}\par}
\begin{document}
\fontsize{9.4}{12.4}\selectfont
{\LARGE\bfseries Plastic plane strain}\par
{\large\color{accent}Boundary conditions and solver settings}\par
{\small Default run: \texttt{coupling/plastic\_plane\_strain.py}; 7 October 2026.}

\topic{Problem and prescribed displacement}
Small-strain, quasi-static J2 plasticity on a $1\,\mathrm{m}\times1\,\mathrm{m}$
square. Body force and all auxiliary sources are zero. Plane strain imposes
zero total $zz$, $xz$ and $yz$ strains; full 3D stress and plastic history are
retained, including nonzero out-of-plane stress.
\[
  \nabla\!\cdot\!\boldsymbol\sigma=0\ \text{in }\Omega,\qquad
  u_x(x,y)=0.004\,\ell\,x,\quad u_y(x,y)=0\ \text{on }\partial\Omega.
\]
Both displacement components are prescribed at every boundary-face centre.
The dimensionless load factor $\ell$ scales the absolute displacement:

\begingroup\renewcommand{\arraystretch}{1.1}
\begin{tabular}{@{}p{0.49\linewidth}p{0.47\linewidth}@{}}
\textbf{Boundary} & \textbf{Prescribed $(u_x,u_y)$}\\\hline
Left: $x=0$ & $(0,0)$\\
Right: $x=1\,\mathrm m$ & $(0.004\,\ell\,\mathrm m,0)$\\
Bottom: $y=0$; top: $y=1\,\mathrm m$ & $(0.004\,\ell\,x,0)$
\end{tabular}\endgroup

This constrains transverse motion: top and bottom also move in $x$ and carry
reactions. No boundary tractions are prescribed. The homogeneous reference has
$\varepsilon_{xx}=0.004\ell$, $\varepsilon_{yy}=\varepsilon_{xy}=0$;
transverse stresses may be nonzero.

\topic{Material and spatial discretization}
$E=210\,\mathrm{GPa}$, $\nu=0.3$, initial yield stress $=250\,\mathrm{MPa}$.
Linear mixed hardening uses $\bar H=1\,\mathrm{GPa}$ and $\theta=0.4$:
the scalar laws are $K(\alpha)=250\,\mathrm{MPa}+(0.4\,\mathrm{GPa})\alpha$
and $H(\alpha)=(0.6\,\mathrm{GPa})\alpha$, where $\alpha$ is accumulated
equivalent plastic strain.

PorePy TPSA uses a fixed $3\times3$ Cartesian grid (9 cells, 36 unknowns):
two displacement components, one rotation stress $r$ and one total pressure $p$
per cell. All $[u,r,p]$ are solved together. Green--Gauss strain reconstruction
feeds one 3D radial return map per cell. Stress corrections use fixed
cell-to-face weights: $1/2$ per neighbour inside; $1$ at the boundary.

\topic{Loading and Newton solve}
Start from zero displacement, stress and history at $\ell=0$. Use 20 prescribed
loading targets $\ell=0.05,0.10,\ldots,1$, then unload to $0.99$ and $0.98$:
22 accepted steps. Peak $\varepsilon_{xx}=0.004$ and right-edge displacement
$=4\,\mathrm{mm}$. Factors label load amplitude; no physical time integration
is used.

Each step starts from the previous converged $[u,r,p]$. Full Newton corrections
use sparse LU (SciPy \texttt{splu}), with at most 20 corrections per step.
The default Jacobian analytically differentiates the 3D return map and TPSA
reconstruction/traction assembly. All three residual blocks must satisfy:
\[
 J_k\,\Delta z_k=-R_k,\quad z_{k+1}=z_k+\Delta z_k,\qquad z=[u,r,p],
\]
\[
 \|R_b^k\|_2\leq a_b+10^{-8}\|R_b^0\|_2,\qquad
 (a_u,a_r,a_p)=(10^{-5},10^{-12},10^{-12}).
\]
These are unweighted Euclidean norms of raw integrated residuals in each
block's SI units; reference norms are fixed at the step's initial evaluation.
There is no damping, line search or adaptive load stepping. Material history
stays fixed during iteration and is accepted after Newton convergence; a failed
solve preserves the preceding state.

Local return-map Newton: at most 20 corrections; tolerance
$=10^{-10}\,\mathrm{Pa}+10^{-12}S$, where $S$ is the larger of the shifted trial
deviatoric-stress norm and $\sqrt{2/3}\,K(\alpha_n)$.
Optional \texttt{--jacobian finite-difference} uses local forward differences
with gradient perturbation $10^{-10}$ (\texttt{--fd-step}), with the same
global assembly.

\topic{Checks and execution}
Every accepted step is checked for equilibrium, agreement with the independent
homogeneous J2 reference, plane strain and elastic unloading. The default run
first detects yield at step 8 ($\ell=0.4$). Newton convergence is printed;
L2 field-error reporting is opt-in (\texttt{--show-l2-errors}).

\texttt{python -m coupling.plastic\_plane\_strain --no-export}

Omit \texttt{--no-export} to write PNGs and VTK/PVD collections under
\texttt{results/}.

\vfill
{\scriptsize Source (repository files): \texttt{coupling/}\{\texttt{plastic\_plane\_strain.py},
\texttt{plane\_strain.py}, \texttt{loading.py}, \texttt{newton.py},
\texttt{residual.py}, \texttt{jacobian.py}, \texttt{transfer.py}\};
\texttt{material\_state.py}, \texttt{scalar\_hardening.py}, \texttt{J2.py}.}
\end{document}
